\part{Topological Lattice Phases}

\section{Gapped Hamiltonians}

\begin{exercise}[A bipartite hopping chain and its bulk spectral gap]
Consider spinless fermions with two distinguished sites $A,B$ per
cell of a one-dimensional lattice. Its one-particle Hamiltonian has
hopping $v>0$ from $A_j$ to $B_j$ and $w>0$ from $B_j$ to $A_{j+1}$,
with Hermitian conjugate hoppings and no on-site terms.
On the Brillouin circle $\T=\R/(2\pi\Z)$ the Fourier transform gives
\[
 H(k)=\begin{pmatrix}0&\overline{q(k)}\\q(k)&0\end{pmatrix},
 \qquad q(k)=v+w e^{\ii k}.
\]
The two entries refer to the physical sublattices. The involution
$\Gamma=+1$ on $A$, $-1$ on $B$ satisfies
$\Gamma H\Gamma=-H$ (chiral symmetry).
A bulk gap at Fermi energy zero is a number $\epsilon>0$ such that
$\Spec H(k)\cap(-\epsilon,\epsilon)=\emptyset$ for every $k$.
\begin{enumerate}[label=(\alph*)]
\item Derive the Fourier formula directly on periodic chains and
prove $H(k)^2=|q(k)|^2\id$. Find the two band energies and their
spectral projections.
\item Prove that a gap exists exactly when $v\ne w$ and that its
half-width is $|v-w|$. At half filling calculate the minimum
particle-hole excitation energy $2|v-w|$.
\item Define the filled-band ground state by occupying the negative
spectral subspace at each allowed momentum. Show that its
one-particle correlation operator is the negative spectral
projection, without choosing eigenvectors.
\end{enumerate}
\end{exercise}

\section{Deformation classes of phases}

\begin{exercise}[Symmetry-preserving deformations and winding number]
Two continuous gapped chiral one-particle Hamiltonian families on
$\T$ are deformation equivalent if they are joined by a continuous
path with a uniform gap and the same chiral grading. Stabilization
allows adding decoupled copies of the constant family
$\left(\begin{smallmatrix}0&1\\1&0\end{smallmatrix}\right)$.
For a smooth nonvanishing scalar $q$ put
\[
 \nu(q)=\frac1{2\pi\ii}\int_{\T}q^{-1}\dd q.
\]
\begin{enumerate}[label=(\alph*)]
\item Lift the phase of $q/|q|$ continuously on $[0,2\pi]$ and
prove that $\nu(q)\in\Z$. Prove its invariance under gapped
deformations and additivity under multiplication.
\item Prove that two scalar loops are homotopic through
nonvanishing loops exactly when their windings agree, by
lifting their quotient's phase. Calculate
$\nu(v+w e^{\ii k})=0$ for $v>w$ and $1$ for $w>v$.
\item For a finite-dimensional graded space with an invertible
off-diagonal map $D(k):V_A\to V_B$, use
$U=D(D^*D)^{-1/2}$ to flatten the Hamiltonian through a
gapped chiral path. For two unitary families $U,U_0$ with
the same domain and codomain define the relative integer
as the winding of $\det(U_0^{-1}U)$, using the determinant
of the section Transfer matrices. Prove stabilization invariance.
\end{enumerate}
\end{exercise}

\section{Generalized cohomology}

\begin{exercise}[Reduced cohomology, suspension, and stable unitary families]
For a based finite cell complex $X$, its reduced suspension is
$\Sigma X=(X\times[0,1])/(X\times\{0,1\}\cup\{*\}\times[0,1])$.
A reduced generalized cohomology theory consists of contravariant
homotopy-invariant graded abelian groups $\widetilde E^n(X)$,
natural suspension isomorphisms
$\widetilde E^{n+1}(\Sigma X)\simeq\widetilde E^n(X)$,
and exact sequences for cell cofibrations $A\hookrightarrow X$
and their quotients $X/A$, with additivity on finite wedges.
\begin{enumerate}[label=(\alph*)]
\item From the cellular cochains of the section Higher gauge
theory construct the relative complex $C^*(X,A;\Z)$ and its
long exact sequence. Compute the cellular cochains of a
suspension to obtain the suspension isomorphism for ordinary
reduced cohomology.
\item A complex vector bundle is a family of vector spaces
locally isomorphic to a product $U\times\C^n$, with linear
fibre identifications. Define $K^0(X)$ as the group completion
of bundles under direct sum and $\widetilde K^0(X)$ as the
kernel of restriction to the base point. Show that a continuous
constant-rank orthogonal projection family on $X$ defines a
bundle, by projecting from a nearby fibre and proving local
invertibility.
\item Define $K^{-1}(X)$ by stable homotopy classes of continuous
unitary endomorphism families, with direct sum as addition.
Show that adjoining identity families does not change the
class, and construct an explicit block-rotation homotopy from
$U\oplus U^{-1}$ to the identity. For $X=S^1$ prove that
determinant winding is an additive homotopy invariant of these
classes. Relate it to the flattened hopping family.
\end{enumerate}
\end{exercise}

\section{Stable homotopy theory}

\begin{exercise}[Suspension spectra and ordinary cohomology]
A sequential spectrum is a sequence of based spaces $E_j$ with
structure maps $\Sigma E_j\to E_{j+1}$. Its stable homotopy
groups are $\pi_n(E)=\varinjlim_j\pi_{n+j}(E_j)$.
The suspension spectrum of $X$ has levels $\Sigma^jX$ and the
canonical structure maps. A loop space $\Omega Y$ consists of
based loops in $Y$; the adjunction
$[\Sigma X,Y]=[X,\Omega Y]$ uses based homotopy classes.
\begin{enumerate}[label=(\alph*)]
\item Prove the suspension-loop adjunction from the two quotient
definitions. Write the transition maps defining $\pi_n(E)$
and show that a compatible levelwise homotopy induces the
same maps on these groups.
\item Use the cocycle models $K(A,j)$ of the section Homotopy
types. Construct the maps
$\Sigma K(A,j)\to K(A,j+1)$ from the suspension of cocycles.
Show that the associated represented groups on finite cell
complexes are ordinary cohomology:
\[
 \widetilde H^n(X;A)=
 \varinjlim_j[\Sigma^jX,K(A,n+j)].
\]
Derive the suspension isomorphism and exactness from the
relative cochain calculation of the preceding section.
\item For $X=S^1$, identify the scalar flattened Hamiltonian
$q/|q|$ with its class in $H^1(S^1;\Z)=\Z$ by phase lifting.
Show that its image under each suspension isomorphism is
the same integer on the suspended fundamental class.
\end{enumerate}
\end{exercise}

\begin{exercise}[Prediction: bulk gap and boundary modes of the hopping chain]
For an open chain of $N$ cells use the hoppings of the section
Gapped Hamiltonians, with no bond from $B_N$ to $A_1$.
Fix $0<v<w$, $r=v/w$, cell spacing $a_0$, chiral symmetry,
and half filling.
\begin{enumerate}[label=(\alph*)]
\item On the semi-infinite chain construct the normalized
zero-energy state supported on sublattice $A$,
$\psi(A_j)=\sqrt{1-r^2}(-r)^{j-1}$, $\psi(B_j)=0$.
Derive its amplitude localization length
$\ell=a_0/\log(w/v)$.
\item On the finite chain normalize the truncated state and
compute its residual norm
\[
 \|H\psi\|=v r^{N-1}
       \sqrt{\frac{1-r^2}{1-r^{2N}}}.
\]
Apply the spectral theorem to prove the existence of a pair
of energies $\pm E_N$ satisfying the same upper bound on
$|E_N|$. Prove that zero itself is absent for $v>0$ by the
invertibility of the triangular hopping map.
\item Connect $v<w$ to $v=0$ through gapped bulk Hamiltonians
and identify the two exactly unpaired boundary sites at that
endpoint. Compare with the gapped deformation of $v>w$ to
$w=0$ and its absence of boundary zero modes.
Relate the difference to the winding class.
\item State the bulk absorption threshold $2(w-v)$ and the
boundary localization length as distinct measurable predictions.
Identify the fixed termination and chiral symmetry as hypotheses;
show directly that an on-site potential at $A_1$ can shift the
boundary energy when that symmetry is relaxed.
\end{enumerate}
\end{exercise}
